\documentclass[12pt, a4paper]{article}

% 引入必要的宏包
\usepackage[UTF8]{ctex}     % 中文支持
\usepackage{amsmath, amssymb, amsthm} % 数学符号和环境
\usepackage{geometry}       % 页面设置
\usepackage{enumitem}       % 列表定制
\usepackage{graphicx}       % 图片支持
\usepackage{fancyhdr}       % 页眉页脚

% 页面边距设置
\geometry{left=2.5cm, right=2.5cm, top=2.5cm, bottom=2.5cm}

% 宏定义：方便排版选择题选项
% 使用方法：\choices{选项A}{选项B}{选项C}{选项D}
\newcommand{\choices}[4]{
    \begin{enumerate}[label=\Alph*., itemsep=0pt, parsep=0pt, topsep=5pt, labelsep=0.5em]
        \item #1
        \item #2
        \item #3
        \item #4
    \end{enumerate}
}
% 横向排列的选项 (2列)
\newcommand{\choicesTwo}[4]{
    \begin{enumerate}[label=\Alph*., itemsep=0pt, parsep=0pt, topsep=5pt, labelsep=0.5em]
        \begin{minipage}{0.45\linewidth} \item #1 \end{minipage}
        \begin{minipage}{0.45\linewidth} \item #2 \end{minipage}
        \begin{minipage}{0.45\linewidth} \item #3 \end{minipage}
        \begin{minipage}{0.45\linewidth} \item #4 \end{minipage}
    \end{enumerate}
}
% 横向排列的选项 (4列)
\newcommand{\choicesFour}[4]{
    \begin{enumerate}[label=\Alph*., itemsep=0pt, parsep=0pt, topsep=5pt, labelsep=0.5em]
        \begin{minipage}{0.22\linewidth} \item #1 \end{minipage}
        \begin{minipage}{0.22\linewidth} \item #2 \end{minipage}
        \begin{minipage}{0.22\linewidth} \item #3 \end{minipage}
        \begin{minipage}{0.22\linewidth} \item #4 \end{minipage}
    \end{enumerate}
}

\title{\textbf{《高等数学》必刷习题——基础版}}
\author{}
\date{}

\begin{document}

\section*{第三章 连续（基础）}

\begin{enumerate}
    \item 函数 $ f(x) $ 在 $\mathbb{R}$ 上连续， $x=2$ 是 $ f(x)=\frac{x^{2}-4}{x-2} $ 的 \underline{\hspace{1cm}}。
    \choicesFour{可去间断点}{跳跃间断点}{无穷间断点}{连续点}

    \item $ x \rightarrow 0^{+} $ 时，下列函数 \underline{\hspace{1cm}} 是无穷小量。
    \choicesFour{$ e^{\frac{1}{x}} $}{$ x\sin\frac{1}{x} $}{$ \ln x $}{$ \frac{1}{x}\sin x $}

    \item 当 $ x \rightarrow 0 $ 时，与 $ x + x^{3} $ 等价的无穷小量为 \underline{\hspace{1cm}}。
    \choicesFour{$ \frac{1}{x} $}{1}{$ x $}{$ x^{2} $}

    \item 函数 $ f(x)=\lim_{n\to\infty}\frac{1+x}{1+x^{2n}} $ 的间断点情况是 \underline{\hspace{1cm}}。
    \choicesTwo
    {不存在间断点}
    {存在间断点 $ x = 1 $}
    {存在间断点 $ x = 0 $}
    {存在间断点 $ x = -1 $}

    \item 设 $ x_{n}=\frac{n}{2}[1+(-1)^{n}] $ ，则 \underline{\hspace{1cm}}。
    \choicesTwo
    {$ \{x_{n}\} $ 有界}
    {$ \{x_{n}\} $ 无界}
    {$ \{x_{n}\} $ 单调增加}
    {$ n\to\infty $ 时， $ \{x_{n}\} $ 为无穷大}

    \item $ x = 0 $ 是函数 $ f(x) = \frac{\tan x}{x} $ 的第 \underline{\hspace{1cm}} 类间断点。

    \item $ f(x)=\begin{cases}x\sin\frac{1}{x}, & x\neq0,\\ k, & x=0\end{cases} $ 在 $x=0$ 处连续,则常数 $k=$ \underline{\hspace{2cm}}。

    \item 已知当 $ x \to 0 $ 时， $ (1 + ax^{2})^{3} - 1 $ 与 $ 1 - \cos x $ 是等价无穷小，则常数 $ a = $ \underline{\hspace{2cm}}。

    \item $ f(x)=\begin{cases}(1-x)^{\frac{1}{x}}, & x<0,\\ 2^{x}+k, & x\geq0\end{cases} $ 在 $ x=0 $ 处连续,则常数 $ k= $ \underline{\hspace{2cm}}。

    \item 函数 $ y=\frac{x}{\tan x} $ 的间断点为 \underline{\hspace{2cm}}。

    \item 已知函数 $ f(x)=\begin{cases}(1+ax)^{\frac{1}{x}}, & x>0\\ 2b-e, & x=0\\ b+\ln(1+x^{2}), & x<0\end{cases} $ 在 $x=0$ 处连续，求实数 $a$ 和 $b$ 的值。

    \item 当 $ a, b $ 为何值时， $ \lim_{x \to 0} \frac{\sin x}{a - e^{x}} (b - \cos x) = 2 $。

    \item $ f(x) = \begin{cases} \frac{\cos x}{x+2}, & x \geq 0 \\ \frac{\sqrt{a} - \sqrt{a - x}}{x}, & x < 0 \quad (a > 0) \end{cases} $
    \begin{enumerate}
        \item 当 $a$ 为何值时， $x=0$ 是 $ f(x) $ 的连续点？
        \item 当 $a$ 为何值时， $x=0$ 是 $ f(x) $ 的间断点？是什么类型的间断点？
    \end{enumerate}

    \item 讨论极限 $ \lim_{x \to 0} \frac{1 - 2^{\frac{1}{x}}}{1 + 2^{\frac{1}{x}}} $ 是否存在？

    \item 研究函数 $ f(x)=\lim_{n\to\infty}\frac{1-x^{2n}}{1+x^{2n}}x $ 的连续性，并画出图形。
\end{enumerate}


\end{document}